% SIAM Article Template
%\providecommand\crefalias[2]{}
\documentclass[review,hidelinks,onefignum,onetabnum]{siamart251216}
%\usepackage{cleveref} 

\usepackage{comment}

\usepackage{amsmath,amssymb,subfigure}      % for subfigures
\usepackage{booktabs,multirow,threeparttable}  %for tables
\usepackage{lipsum}
\usepackage{amsfonts}
\usepackage{graphicx}
\usepackage{epstopdf}
\usepackage{algorithmic}
\usepackage{mathrsfs} 
\usepackage{multirow}

\ifpdf
  \DeclareGraphicsExtensions{.eps,.pdf,.png,.jpg}
\else
  \DeclareGraphicsExtensions{.eps}
\fi

% Add a serial/Oxford comma by default.
\newcommand{\creflastconjunction}{, and~}

% Used for creating new theorem and remark environments
\newsiamremark{remark}{Remark}
\newsiamremark{Assumption}{Assumption}
\newsiamremark{assumption}{Assumption}
\crefname{Assumption}{Assumption}{Hypotheses}
\newsiamthm{claim}{Claim}


\usepackage{amsopn}
\DeclareMathOperator{\diag}{diag}

\newcommand{\ud}{\mathrm{d}}
\newtheorem{example}{\hspace{6mm}Example}[section]
\renewcommand{\div}{\operatorname{div}}
\graphicspath{{./figure/}{../figure/}}

\newcommand{\dx}{\,{\rm d}x}
\newcommand{\ds}{\,{\rm d}s}
\newcommand{\dd}{\,{\rm d}}
\newcommand{\Oplus}{\ensuremath{\vcenter{\hbox{\scalebox{1.5}{$\oplus$}}}}}

\newcommand{\tr}{\operatorname{tr}}
\newcommand{\alt}{\operatorname{Alt}}
\newcommand{\dev}{\operatorname{dev}}
\newcommand{\sym}{\operatorname{sym}}
\newcommand{\skw}{\operatorname{skw}}
\newcommand{\spn}{\operatorname{spn}}
\newcommand{\mspn}{\operatorname{mspn}}
\newcommand{\vspn}{\operatorname{vspn}}
\newcommand{\mskw}{\operatorname{mskw}}
\newcommand{\vskw}{\operatorname{vskw}}
\newcommand{\defm}{\operatorname{def}}

%%%% Comments of authors %%%%%
\newcommand{\LC}[1]{\textcolor{red}{#1}}
\newcommand{\XH}[1]{\textcolor{blue}{#1}}
\newcommand{\RW}[1]{\textcolor{cyan}{#1}}
\newcommand{\RZ}[1]{\textcolor{brown}{#1}}
% Substantive proofreading changes; use #1 alone for a clean copy.
\newcommand{\AI}[1]{{\color{blue}#1}}



\title{A Stabilizer-Free Weak Virtual Element Method for the
Stokes Problem on Polytopal Meshes\thanks{Submitted to the editors DATE.
\funding{The first author was supported by ...}}}

% Authors: full names plus addresses.
\author{Eun-Jae Park\thanks{Department of Computational Science and Engineering, 
           Yonsei University, 
           Seoul 03722, Korea
  (\email{ejpark@yonsei.ac.kr}).}
\and Feng Wang\thanks{Jiangsu Key Laboratory for NSLSCS, 
School of Mathematical Sciences, Nanjing Normal University,
Nanjing 210023, China  
  (\email{fengwang@live.cn}).}
\and Ruishu Wang\thanks{School of Mathematics,
Jilin University, 
Changchun, Jilin 130012, China
  (\email{wangrs\_math@jlu.edu.cn}).}
}

\begin{document}



\maketitle




%\maketitle
%
%% REQUIRED
%\begin{abstract}
%This paper develops a novel 
%\end{abstract}
%
%% REQUIRED
%\begin{keywords}
% Error estimate
%\end{keywords}

% REQUIRED
%\begin{MSCcodes}
%65N30, 65N15, 65N12, 74B05
%\end{MSCcodes}
%65N30: Finite elements, Rayleigh-Ritz and Galerkin methods, finite methods
%65N15: Error bounds
%65N12: Stability and convergence of numerical methods
%74B05: Classical linear elasticity

\section{Introduction}\label{sec:introduction}

\AI{A discretization of incompressible Stokes flow must preserve mass and
control the velocity through its symmetric gradient. On polygonal and
polyhedral meshes, these requirements call for discrete spaces with suitable
normal and tangential traces. The virtual element method provides such spaces
on general polytopes. Its degrees of freedom determine the polynomial
projections needed to assemble the discrete problem without evaluating the
virtual basis functions \cite{BBCMMR2013}. In standard formulations, a
stabilization term controls the part of a virtual function that its polynomial
projection does not capture.

The $H(\operatorname{div})$-conforming virtual elements of
Beir\~ao da Veiga, Brezzi, Marini, and Russo \cite{BBMR2016} provide continuous
normal traces and polynomial divergence. Chen and Wang \cite{ChenWang2019}
combined these elements with tangential face unknowns and a weak symmetric
gradient to obtain a divergence-free weak virtual element method for the
Stokes problem. Their formulation includes a stabilization term that controls
the boundary mismatch between the projected cell velocity and its traces.
This framework gives a natural starting point for a method whose weak strain
alone controls that mismatch.

Stabilizer-free methods achieve stability through the choice of the local
reconstruction space. Mu, Ye, and Zhang \cite{MuYeZhang2021} developed a
stabilizer-free weak Galerkin method for the Stokes equations on polytopal
meshes. For second-order elliptic problems, Chen, Huang, and Wei
\cite{ChenHuangWei2024} constructed virtual element methods without extrinsic
stabilization by using stable gradient projections onto macro finite element
spaces. These results motivate local reconstructions that retain enough
information to control the discrete energy. For the symmetric-gradient
formulation of Stokes flow, the reconstruction must also account for rigid
motions and the coupling of normal and tangential velocity traces.

Here we develop a stabilizer-free weak virtual element method in two and
three dimensions. The velocity consists of an $H(\operatorname{div})$ virtual
field and a tangential polynomial field on the primal mesh faces. The pressure
is a discontinuous polynomial of degree $k-1$. We reconstruct the weak strain
in a symmetric tensor space that is polynomial on the cones joining each
element's star center to its faces. The strain is computed locally from the
velocity degrees of freedom. The tensor space provides independent face
moments and volume moments that control both the symmetric gradient of the
projected velocity and the boundary mismatch. The discrete viscous energy
is the squared $L^2$ norm of this weak strain.

Under the stated mesh regularity and a uniform spanning condition on the
face-tangential symmetric tensors, we prove norm equivalence, coercivity,
and the velocity--pressure inf-sup condition. The discrete velocity is exactly
divergence-free on every primal element. For sufficiently smooth solutions,
we obtain order-$h^k$ error bounds in the discrete velocity norm and the
$L^2$ pressure norm. The spanning condition is an explicit assumption of the
analysis. The numerical section reports two-dimensional tests and identifies
the rectangular-mesh results as tests outside this assumption.

Section~\ref{sec:spaces} introduces the meshes, velocity spaces, and
projections. Section~\ref{sec:strain} constructs the weak strain and proves
its stability. Section~\ref{sec:scheme} presents the discrete problem and
error estimates. Section~\ref{sec:numerics} gives the numerical results.}

\section{VEM spaces}\label{sec:spaces}

Let \(\Omega \subset \mathbb{R}^d\), \(d = 2, 3\) be a bounded and simply connected polytopal domain. Given an
external force field \(\boldsymbol{f} \in \boldsymbol{L}^2(\Omega)\), we consider the steady-state incompressible Stokes problem:
\begin{equation}\label{eq:stokes}
\left\{
\begin{aligned}
-\operatorname{div}(\varepsilon(\boldsymbol{u})) + \nabla p &= \boldsymbol{f} \quad \text{in } \Omega, \\
\operatorname{div} \boldsymbol{u} &= 0 \quad \text{in } \Omega, \\
\boldsymbol{u} &= \boldsymbol{0} \quad \text{on } \partial\Omega.
\end{aligned}
\right.
\end{equation}


\AI{Here $\varepsilon(\boldsymbol u):=(\nabla\boldsymbol u+
\nabla\boldsymbol u^{\mathsf T})/2$ is the symmetric strain, and the pressure
is normalized by $\int_\Omega p\,\mathrm dx=0$. We use $(\cdot,\cdot)_K$ and
$\langle\cdot,\cdot\rangle_{\partial K}$ for the volume and boundary $L^2$
inner products, and $\|\cdot\|_K$ and $\|\cdot\|_{\partial K}$ for their norms.
}

\subsection{Primal and dual meshes}
\AI{Let $\mathcal K_h$ be a conforming polytopal mesh of $\Omega$.
Each element $K$ is star-shaped with respect to a ball of radius at least
$\rho h_K$, where $h_K=\operatorname{diam}K$ and $\rho>0$ is independent of
$h:=\max_K h_K$. Let $\boldsymbol x_K$ be the center of such a ball.
We assume a uniformly bounded number of faces per element. In three dimensions,
each face is also uniformly star-shaped. The cones
$T=\{\boldsymbol x_K+t(\boldsymbol y-\boldsymbol x_K):\boldsymbol y\in F,\ 0\leq t\leq1\}$ form a refinement $\mathcal T_h$;
they are triangles in two dimensions and pyramids in three dimensions.
We assume that these cones admit a conforming, uniformly shape-regular
simplicial refinement with simplex diameters comparable to $h_K$.
In particular, $h_F:=\operatorname{diam}F$ is comparable to $h_K$.
These assumptions are used in the trace, inverse, and scaling estimates below.}

Let $\mathcal F_h$ and $\mathcal F_h^T$ be the face sets of $\mathcal K_h$
and $\mathcal T_h$, respectively. Their interior faces are denoted by
$\mathring{\mathcal F}_h^K$ and $\mathring{\mathcal F}_h^T$.
Let $\mathcal T_h(K)$ be the set of cones in $K$.

Let $F \in \partial K$ be a face with unit normal $\boldsymbol{n}_F$ and orthonormal tangent vectors $\{\boldsymbol{t}_i^{F}\}_{i=1}^{d-1}$. For any $\boldsymbol{v}\in\mathbb{R}^d$, the tangential part on $F$ is
$$
\Pi_F \boldsymbol{v}=(\boldsymbol{I}-\boldsymbol{n}_F\otimes\boldsymbol{n}_F)\boldsymbol{v}
=\sum_{i=1}^{d-1}(\boldsymbol{v}\cdot\boldsymbol{t}_i^{F})\boldsymbol{t}_i^{F}.
$$

With $\sym A=(A+A^{\mathsf T})/2$ and $\skw A=(A-A^{\mathsf T})/2$, define
\begin{align*}
\mathscr T^F&:=\textrm{span}\{\boldsymbol{t}_1^F, \ldots, \boldsymbol{t}_{d-1}^F\}, \\
\mathscr T^F(\mathbb S)&:=\textrm{span}\{{\rm sym}(\boldsymbol{t}_i^F\otimes \boldsymbol{t}_j^F) \;\;\textrm{ for } 1\leq i\leq j\leq d-1\}, \\
\mathscr N^F(\mathbb S) &:= \sym(\mathscr T^F\otimes\boldsymbol{n}_F)\oplus {\rm span}\{\boldsymbol{n}_F\otimes\boldsymbol{n}_F\},\\
\mathscr T^F(\mathbb K)&:=\textrm{span}\{{\rm skw}(\boldsymbol{t}_i^F\otimes \boldsymbol{t}_j^F) \;\;\textrm{ for } 1\leq i< j\leq d-1\}.
\end{align*}
These subspaces give the following $\boldsymbol{t}$-$\boldsymbol{n}$ decompositions
\begin{align}\label{eq:decS}
\mathbb S&= \mathscr T^F(\mathbb S)\oplus \mathscr N^F(\mathbb S), \quad 
\\
\mathbb K&= \mathscr T^F(\mathbb K)\oplus \mathscr N^F(\mathbb K), \quad \mathscr N^F(\mathbb K)={\rm skw}(\mathscr T^F\otimes\boldsymbol{n}_F).
\end{align}


\AI{Fix an integer $k\geq1$.} For \(K \in \mathcal{K}_h\), let \(\mathbb{P}_k(K)\) be the space of polynomials on \(K\) of degree at most \(k\). Let \(\mathbb{P}_k(K; \mathbb{X}) = \mathbb{P}_k(K) \otimes \mathbb{X}\), where \(\mathbb{X}\) is \(\mathbb{R}^d\) for vectors, \(\mathbb{S}\) for symmetric matrices, and \(\mathbb{K}\) for skew-symmetric matrices.

\AI{We set $\mathbb P_j=\{0\}$ for $j<0$, and
$\mathbf{RM}(K)=\{\boldsymbol a+\boldsymbol B\boldsymbol x:
\boldsymbol a\in\mathbb R^d,\ \boldsymbol B\in\mathbb K\}$.}

\subsection{VEM spaces}

\AI{The local velocity space and its degrees of freedom are the
$H(\operatorname{div})$ virtual elements of \cite{BBMR2016}.}
On each element \(K\in\mathcal{K}_h\),  we introduce a local space
\begin{equation*}
\boldsymbol{V}_K = \left\{ 
\begin{aligned}
& \boldsymbol{v} \in \boldsymbol{H}(\operatorname{div}, K) \cap \boldsymbol{H}(\operatorname{curl}, K) : \boldsymbol{v} \cdot \boldsymbol{n} \in \mathbb{P}_k(F) \ \forall F \subset \partial K, \\
& \operatorname{div} \boldsymbol{v} \in \mathbb{P}_{k-1}(K), \ \operatorname{curl} \boldsymbol{v} \in \operatorname{curl} \mathbb{P}_k(K;\mathbb{R}^d)
\end{aligned}
\right\},
\end{equation*}

The global virtual element space is
\begin{equation*}
\boldsymbol{V}_h^{\operatorname{div}} := \{ \boldsymbol{v} \in \boldsymbol{H}(\operatorname{div}, \Omega), \boldsymbol{v}|_K \in \boldsymbol{V}_K \ \forall K \in \mathcal{K}_h, \boldsymbol{v} \cdot \boldsymbol{n} = 0 \text{ on } \partial \Omega \} .
\end{equation*}


The degrees of freedom (DoFs) for $\boldsymbol V_K$ are
\begin{align*}
\text{Type I} \qquad & \int_F \boldsymbol{v} \cdot \boldsymbol{n} q_k \, \mathrm{d}s \qquad \forall q_k \in \mathbb{P}_k(F), F \subset \partial K, \\
\text{Type II} \qquad & \int_K \boldsymbol{v} \cdot \boldsymbol{q}_{k-2} \, \mathrm{d}\boldsymbol{x} \qquad \forall \boldsymbol{q}_{k-2} \in \mathbf{G}_{k-2}(K), \\
\text{Type III} \qquad & \int_K \boldsymbol{v} \cdot \boldsymbol{q}_k \, \mathrm{d}\boldsymbol{x} \qquad \forall \boldsymbol{q}_k \in \AI{\mathbf G_k(K)^\perp},
\end{align*}
where $\mathbf G_k(K)=\nabla\mathbb P_{k+1}(K)$.
\AI{Here $\mathbf G_k(K)^\perp$ is its $L^2(K)$-orthogonal complement in
$\mathbb P_k(K;\mathbb R^d)$. Moments are taken against a basis of each
listed space. These DoFs are unisolvent \cite{BBMR2016}.}

On each face \(F \in \mathcal{F}_h\), we introduce a space 
$$
\boldsymbol{V}_F := \mathbb{P}_{k-1}(F;\mathscr T^F) + \mathbf{RM}(F),
$$
 where \(\mathbf{RM}(F)\) denotes the space of rigid motions, i.e., \(\mathbf{RM}(F) :=\Pi_F \mathbf{RM}(K)= \mathbb{P}_0(F; \mathscr T^F) \oplus \mathbb{P}_0(F; \mathscr T^F(\mathbb K)) \boldsymbol{x}\). The DoFs for functions in \(\boldsymbol{V}_F\) are
\[
\int_F \boldsymbol{v} \cdot \boldsymbol{q} \, \mathrm{d}s \quad \forall \boldsymbol{q} \in \mathbb{P}_{k-1}(F;\mathscr T^F) + \mathbf{RM}(F).
\]

The tangential face space is
\[
\boldsymbol{V}_h^t := \{ \boldsymbol{v}^t, \boldsymbol{v}^t|_F \in \boldsymbol{V}_F, F \in \mathcal{F}_h, \boldsymbol{v}^t = \boldsymbol{0} \text{ on } \partial \Omega \}.
\]

The weak virtual element space is
\[
\boldsymbol{V}_h := \left\{ \{ \boldsymbol{v}_h^{\operatorname{div}}, \boldsymbol{v}_h^t \}, \boldsymbol{v}_h^{\operatorname{div}} \in \boldsymbol{V}_h^{\operatorname{div}}, \boldsymbol{v}_h^t \in \boldsymbol{V}_h^t \right\}.
\]

\subsection{Projection operator}

\AI{For $\boldsymbol v\in\boldsymbol L^2(K)$,} let \(\boldsymbol{\Pi}_K^o \boldsymbol{v} \in \mathbb{P}_k(K;\mathbb{R}^d)\) satisfy
\[
\int_K \boldsymbol{\Pi}_K^o \boldsymbol{v} \cdot \boldsymbol{q}\,\mathrm dx = \int_K \boldsymbol{v} \cdot \boldsymbol{q}\,\mathrm dx \quad \forall \boldsymbol{q} \in \mathbb{P}_k(K;\mathbb{R}^d).
\]

For $\boldsymbol v\in\boldsymbol V_K$, Type III DoFs give the moments in $\mathbf G_k(K)^\perp$. The remaining moments follow by integration by parts: for $p_{k+1}\in\mathbb P_{k+1}(K)$,
\[
\int_K \boldsymbol{v} \cdot \nabla p_{k+1} \, \mathrm{d}\boldsymbol{x} = - \int_K \operatorname{div} \boldsymbol{v} \, p_{k+1} \, \mathrm{d}\boldsymbol{x} + \int_{\partial K} \boldsymbol{v} \cdot \boldsymbol{n} \, p_{k+1} \, \mathrm{d}s.
\]

\AI{The normal trace is determined by Type I DoFs. The divergence is determined
by Type I and II DoFs through
\[
(\operatorname{div}\boldsymbol v,q)_K
=-(\boldsymbol v,\nabla q)_K+\langle\boldsymbol v\cdot\boldsymbol n,q\rangle_{\partial K},
\qquad q\in\mathbb P_{k-1}(K).
\]
Thus every term above is computable \cite{BBMR2016}. We write
$(\boldsymbol\Pi_h^o\boldsymbol v)|_K=\boldsymbol\Pi_K^o(\boldsymbol v|_K)$.}



\section{Weak strain}\label{sec:strain}

\subsection{Weak strain space}
\AI{Let $T\in\mathcal T_h(K)$ have base face $F$.
The local weak strain space is
\begin{equation}\label{eq:strainspace}
\begin{aligned}
\Sigma_k(T)={}&\mathbb P_k(T)\boldsymbol n_F\otimes\boldsymbol n_F\\
&\oplus\sym\!\left(\boldsymbol n_F\otimes
  [\mathbb P_{k-1}(T;\mathscr T^F)+\mathbf{RM}(F)]\right)\\
&\oplus\mathbb P_{k-1}(T;\mathscr T^F(\mathbb S)).
\end{aligned}
\end{equation}
Face polynomials, including $\mathbf{RM}(F)$, are extended constantly along
lines normal to $F$. For $k=1$, the middle block is
$\sym(\boldsymbol n_F\otimes\mathbf{RM}(F))$; for $k>1$, it is
$\sym(\boldsymbol n_F\otimes\mathbb P_{k-1}(T;\mathscr T^F))$.
Set
\[
V_k(F):=\mathbb P_k(F)\boldsymbol n_F\oplus\boldsymbol V_F.
\]
The normal trace $\tr_F\boldsymbol\sigma:=\boldsymbol\sigma\boldsymbol n_F|_F$
satisfies $\tr_F\Sigma_k(T)=V_k(F)$. In particular,
$\mathbb P_{k-1}(K;\mathbb S)$ is contained in the piecewise strain space on $K$.}

We define the local symmetric tensor space on $K$ as the discontinuous polynomial space
\begin{align}\label{eq:202411132}
\Sigma_h^{-1}(K)&=\Oplus_{i=1}^n\chi_{T_i}\Sigma_k(T_i)=\{\boldsymbol{\tau}\in L^2(K;\mathbb S): \boldsymbol{\tau}|_T\in \Sigma_k(T) \textrm{ for }T\in\mathcal{T}_h(K)\},
\end{align}
where $\chi_{T_i}$ is the characteristic function of $T_i$. We  define the bubble space
\begin{equation}
\begin{aligned}
\mathbb{B}_k(\div,K;\mathbb{S})=&\{\boldsymbol{\tau}\in \Sigma_h^{-1}(K): \boldsymbol{\tau}\boldsymbol{n}|_{\partial K}=\boldsymbol{0}\}\\
=&\Oplus_{i=1}^n\lambda_{\boldsymbol{x}_K}\chi_{T_i}(\mathbb{P}_{k-1}(T_i)\boldsymbol{n}_i\otimes\boldsymbol{n}_i) \\
\oplus
&\Oplus_{i=1}^n\lambda_{\boldsymbol{x}_K}\chi_{T_i}\sym(\mathbb{P}_{k-2}(T_i; \mathscr T^{F_i})\otimes\boldsymbol{n}_i) \\
\oplus 
&\Oplus_{i=1}^n\chi_{T_i}\mathbb{P}_{k-1}(T_i;\mathscr T^{F_i}(\mathbb S)).
\end{aligned}
\end{equation}

\AI{Here $F_i$ is the base face of $T_i$, $\boldsymbol n_i$ is its outward normal,
and $\lambda_{\boldsymbol x_K}$ is the affine function on $T_i$ that equals
one at $\boldsymbol x_K$ and zero on $F_i$. The bubble space imposes no
continuity across the internal faces of $K$.}

\begin{assumption}\label{meshassumption}
There exist $d(d+1)/2$ symmetric tensors $b_\ell$ such that
\begin{equation}
\mathbb{S}=\text{span}\left\{ b_{\ell}, \ell=1,\ldots,d(d+1)/2\right\},\label{equ_assumption_ddpolytope}
\end{equation}
where $\mathscr T^{F_i}(\mathbb S)=\textrm{span}\{ b_i^{m}, m=1,\ldots,d(d-1)/2\}$ for $F_i\in \partial K$, and 
\begin{equation}
\left\{ b_\ell\right\} \subset \left\{ b_i^m, i=1,\ldots,n, m=1,\ldots,d(d-1)/2\right\}.\label{equ_edge_vector_property}
\end{equation}
\AI{On every element, the tensors $b_\ell$ can be chosen with
$|b_\ell|=1$ and
\[
\sum_{\ell=1}^{d(d+1)/2}c_\ell^2
\lesssim\left|\sum_{\ell=1}^{d(d+1)/2}c_\ell b_\ell\right|^2
\qquad\text{for all real coefficients }c_\ell,
\]
with a constant independent of $h$ and $K$. Here $|\cdot|$ is the Frobenius norm.}
\end{assumption}
\AI{All subsequent stability and error estimates assume
Assumption~\ref{meshassumption}. Constants may depend on $k$, $\Omega$, the
mesh regularity, and this tensor-basis bound, but not on $h$.}

\AI{For each selected $b_\ell$, let $T_\ell$ be a cone whose base face
has $b_\ell\in\mathscr T^{F_i}(\mathbb S)$, and define
\[
\mathcal E_K:=\bigoplus_{\ell=1}^{d(d+1)/2}
\chi_{T_\ell}\mathbb P_{k-1}(T_\ell)b_\ell
\subset\mathbb B_k(\operatorname{div},K;\mathbb S).
\]
In the following lemma, $\mathbb B_k\ominus\mathcal E_K$ denotes the
$L^2(K)$-orthogonal complement of $\mathcal E_K$ within the bubble space.}

\begin{lemma}\label{lemma_DOF}
For $K\in\mathcal{K}_h$, the space $\Sigma_h^{-1}(K)$ is uniquely determined by the following DoFs:
\begin{subequations}
\label{SigmaKDoFs}
\begin{align}
\label{SigmaKDoFs1}
\int_{F}\boldsymbol{\sigma}\boldsymbol{n}\cdot \boldsymbol{q}\ds,&\quad \boldsymbol{q}\in   V_{k}(F), F\in \partial K , \\
\label{SigmaKDoFs2}
\int_K\boldsymbol{\sigma}:\boldsymbol{\tau} \dx,& \quad  \AI{\boldsymbol\tau\in\mathbb B_k(\div,K;\mathbb S)\ominus\mathcal E_K}, \\
\label{SigmaKDoFs3}
\int_K\boldsymbol{\sigma}:\boldsymbol{\tau} \dx,& \quad  \boldsymbol{\tau}\in \mathbb{P}_{k-1}(K;\mathbb{S}),
\end{align}
\end{subequations}
where moments are taken against bases of the listed spaces.
\end{lemma}
\AI{\begin{proof}
Suppose all moments vanish. Since the normal trace belongs to $V_k(F)$,
the face moments imply $\boldsymbol\sigma\boldsymbol n=0$ on $\partial K$.
Thus $\boldsymbol\sigma$ is a bubble. The second group of moments then gives
$\boldsymbol\sigma\in\mathcal E_K$, so
$\boldsymbol\sigma=\sum_\ell\chi_{T_\ell}p_\ell b_\ell$ with
$p_\ell\in\mathbb P_{k-1}(T_\ell)$.
Let $b_\ell^*$ be the dual tensor basis, satisfying
$b_j:b_\ell^*=\delta_{j\ell}$. Extend $p_\ell$ as a polynomial to $K$ and
choose $\boldsymbol\tau=p_\ell b_\ell^*$ in \eqref{SigmaKDoFs3}. Then
$\|p_\ell\|_{T_\ell}^2=0$ for every $\ell$, so $\boldsymbol\sigma=0$.
The trace map is onto $\bigoplus_{F\subset\partial K}V_k(F)$ and has the
bubble space as its kernel. Also,
$\dim\mathcal E_K=\dim\mathbb P_{k-1}(K;\mathbb S)$.
The number of moments therefore equals $\dim\Sigma_h^{-1}(K)$.
This proves unisolvence.
\end{proof}}

Let
\[
\Sigma_h^{-1} = \{\boldsymbol{\sigma} \in L^2(\Omega;\mathbb{S}) : \boldsymbol{\sigma}|_T \in \Sigma_k(T), T \in \mathcal{T}_h\}.
\]

\subsection{Weak strain}
For any \(\boldsymbol{v} = \{ \boldsymbol{v}^{\operatorname{div}}, \boldsymbol{v}^t \} \in \boldsymbol{V}_h\), 
we define \(\varepsilon^w_h(\boldsymbol{v}) \in \Sigma_h^{-1}\) satisfying \( \varepsilon^w_h(\boldsymbol{v})|_K=\varepsilon_K^w(\boldsymbol{v})\) and
\begin{equation}\label{equ:weakstrain1}
\begin{aligned}
(\varepsilon_K^w(\boldsymbol v),\boldsymbol\sigma)_K
={}&(\varepsilon(\boldsymbol\Pi_K^o\boldsymbol v^{\operatorname{div}}),\boldsymbol\sigma)_K\\
&+\langle(\boldsymbol v^{\operatorname{div}}-
 \boldsymbol\Pi_K^o\boldsymbol v^{\operatorname{div}})\cdot\boldsymbol n,
 (\boldsymbol\sigma\boldsymbol n)\cdot\boldsymbol n\rangle_{\partial K}\\
&+\langle\boldsymbol v^t-\Pi_F\boldsymbol\Pi_K^o\boldsymbol v^{\operatorname{div}},
 \Pi_F(\boldsymbol\sigma\boldsymbol n)\rangle_{\partial K}
\end{aligned}
\end{equation}
for all $\boldsymbol\sigma\in\Sigma_h^{-1}(K)$, with $\Pi_F$ applied facewise.
\AI{Since $\boldsymbol\sigma$ is discontinuous inside $K$, integration by
parts must include the internal jumps of $\boldsymbol\sigma\boldsymbol n$.
Dropping these terms does not give an equivalent definition.}

The normal and tangential boundary jumps are
\[
\begin{aligned}
\mathcal J_{\partial K}^{\boldsymbol n}(\boldsymbol v)
&:=((\boldsymbol v^{\operatorname{div}}-
\boldsymbol\Pi_K^o\boldsymbol v^{\operatorname{div}})\cdot\boldsymbol n)\boldsymbol n,\\
\mathcal J_{\partial K}^{\boldsymbol\tau}(\boldsymbol v)
&:=\boldsymbol v^t-\Pi_F\boldsymbol\Pi_K^o\boldsymbol v^{\operatorname{div}}.
\end{aligned}
\]
Their sum is $\mathcal J_{\partial K}(\boldsymbol v):=
\mathcal J_{\partial K}^{\boldsymbol n}(\boldsymbol v)+
\mathcal J_{\partial K}^{\boldsymbol\tau}(\boldsymbol v)$.

 On the face \(F \subset \partial K\), we define \((\boldsymbol{\Pi}_{\partial K}^t \boldsymbol{v})|_F \in \boldsymbol{V}_F\) as
\[
\int_F \boldsymbol{\Pi}_{\partial K}^t \boldsymbol{v} \cdot \boldsymbol{q}\,\mathrm ds = \int_F \boldsymbol{v}|_K \cdot \boldsymbol{q}\,\mathrm ds \quad \forall \boldsymbol{q} \in \boldsymbol{V}_F.
\]
Define the notation \(\boldsymbol{\Pi}_{\partial K}\) as
\[
\boldsymbol{\Pi}_{\partial K} \boldsymbol{v} := (\boldsymbol{v} \cdot \boldsymbol{n})\boldsymbol{n} + \boldsymbol{\Pi}_{\partial K}^t \boldsymbol{v}.
\]

Since $\boldsymbol\sigma\boldsymbol n|_F\in V_k(F)$, \eqref{equ:weakstrain1} becomes
\[
(\varepsilon_K^w(\boldsymbol{v}), \boldsymbol{\sigma})_K = (\varepsilon(\boldsymbol{\Pi}_K^o \boldsymbol{v}^{\operatorname{div}}), \boldsymbol{\sigma})_K + \langle \boldsymbol{\Pi}_{\partial K} \mathcal{J}_{\partial K}(\boldsymbol{v}), \boldsymbol{\sigma}\boldsymbol{n} \rangle_{\partial K}.
\]


For $\boldsymbol v_h\in\boldsymbol V_h$, define the mesh-dependent seminorm
\[
\|\boldsymbol{v}_h\|_{1,h}^2 := \sum_{K \in \mathcal{K}_h} \|\boldsymbol{v}_h\|_{1,h,K}^2,
\]
where the seminorm on each element is defined as
\[
\AI{\|\boldsymbol{v}_h\|_{1,h,K}^2 := \|\varepsilon (\boldsymbol{\Pi}_K^o \boldsymbol{v}_h^{\operatorname{div}})\|_K^2 + h_K^{-1}\|\boldsymbol{\Pi}_{\partial K}\mathcal{J}_{\partial K}(\boldsymbol{v}_h)\|_{\partial K}^2.}
\]


\begin{lemma}\label{lem:norm}
The mesh-dependent quantity
\(
\|\boldsymbol{v}_h\|_{1,h}
\)
defines a norm on the virtual element space \(\boldsymbol{V}_h\).
\end{lemma}

\AI{\begin{proof}
If $\|\boldsymbol v_h\|_{1,h}=0$, then on each $K$,
$\boldsymbol\Pi_K^o\boldsymbol v_h^{\operatorname{div}}=\boldsymbol r_K
\in\mathbf{RM}(K)$, and
\[
\boldsymbol v_h^{\operatorname{div}}\cdot\boldsymbol n
=\boldsymbol r_K\cdot\boldsymbol n,
\qquad
\boldsymbol v_h^t=\Pi_F\boldsymbol r_K\quad\text{on each }F\subset\partial K.
\]
The second identity uses $\Pi_F\boldsymbol r_K\in\mathbf{RM}(F)$.
The normal component of $\boldsymbol v_h^{\operatorname{div}}$ and the
single-valued tangential field $\boldsymbol v_h^t$ show that neighboring
rigid motions agree on their common face. A rigid motion that vanishes on
a face is zero, so connectedness gives one global rigid motion $\boldsymbol r$.
The boundary conditions give both $\boldsymbol r\cdot\boldsymbol n=0$ and
$\Pi_F\boldsymbol r=0$ on boundary faces. Hence $\boldsymbol r=0$.
The normal face moments and all polynomial volume moments of
$\boldsymbol v_h^{\operatorname{div}}$ now vanish. Unisolvence of
$\boldsymbol V_K$ gives $\boldsymbol v_h^{\operatorname{div}}=0$,
and the tangential identity gives $\boldsymbol v_h^t=0$.
\end{proof}}

\AI{We also need control of the full jump. For
$\boldsymbol w=\boldsymbol\Pi_K^o\boldsymbol v_h^{\operatorname{div}}$,
Korn's inequality modulo rigid motions gives a $\boldsymbol r_K\in\mathbf{RM}(K)$
such that
\[
h_K^{-1}\|\boldsymbol w-\boldsymbol r_K\|_K
+\|\nabla(\boldsymbol w-\boldsymbol r_K)\|_K
\lesssim\|\varepsilon(\boldsymbol w)\|_K.
\]
Since $\boldsymbol\Pi_{\partial K}^t$ reproduces $\Pi_F\boldsymbol r_K$,
the trace inequality yields
\[
h_K^{-1/2}\|\Pi_F\boldsymbol w-
\boldsymbol\Pi_{\partial K}^t\boldsymbol w\|_{\partial K}
\lesssim\|\varepsilon(\boldsymbol w)\|_K.
\]
The normal part of the jump is unchanged by $\boldsymbol\Pi_{\partial K}$.
The unprojected tangential part has the same norm as the difference above. Thus
\begin{equation}\label{eq:fulljump}
h_K^{-1/2}\|\mathcal J_{\partial K}(\boldsymbol v_h)\|_{\partial K}
\lesssim\|\boldsymbol v_h\|_{1,h,K}.
\end{equation}
The jump of $\boldsymbol w$ across a primal face is the difference of its
two full boundary jumps; on a boundary face, $\boldsymbol w$ is minus its
full boundary jump. Applying the broken Korn inequality \cite{Brenner2004}
on the simplicial refinement therefore gives the global equivalence
\begin{equation}\label{eq:globalKorn}
\|\boldsymbol v_h\|_{1,h}^2\simeq
\sum_{K\in\mathcal K_h}\left(
\|\nabla\boldsymbol\Pi_K^o\boldsymbol v_h^{\operatorname{div}}\|_K^2
+h_K^{-1}\|\boldsymbol\Pi_{\partial K}\mathcal J_{\partial K}(\boldsymbol v_h)\|_{\partial K}^2
\right).
\end{equation}
This equivalence is global; it does not hold on an individual element.}

\begin{lemma} \label{lem:discinf-sup}
We have the discrete inf-sup condition
\begin{equation}\label{eq:discinfsup}
\|\boldsymbol{v}\|_{1,h}\lesssim \sup_{0\ne\boldsymbol{\sigma}_h\in\Sigma_h^{-1}}
\frac{\sum_{K\in\mathcal{K}_h}(\varepsilon_K^w(\boldsymbol{v}), \boldsymbol{\sigma}_h)_K}{\|\boldsymbol{\sigma}_h\|_{0,h}}\quad\forall~\boldsymbol{v}\in\boldsymbol V_h.
\end{equation}
\end{lemma}
\AI{\begin{proof}
Here $\|\boldsymbol\sigma\|_{0,h}^2:=\sum_K\|\boldsymbol\sigma\|_K^2$.
For fixed $\boldsymbol v\in\boldsymbol V_h$, use Lemma~\ref{lemma_DOF}
to choose $\boldsymbol\sigma|_K\in\Sigma_h^{-1}(K)$ with
\[
\begin{aligned}
(\boldsymbol\sigma,\boldsymbol\tau)_K
&=(\varepsilon(\boldsymbol\Pi_K^o\boldsymbol v^{\operatorname{div}}),\boldsymbol\tau)_K
&&\forall\boldsymbol\tau\in\mathbb P_{k-1}(K;\mathbb S),\\
\langle\boldsymbol\sigma\boldsymbol n,\boldsymbol q\rangle_F
&=\langle h_K^{-1}\boldsymbol\Pi_{\partial K}\mathcal J_{\partial K}(\boldsymbol v),\boldsymbol q\rangle_F
&&\forall\boldsymbol q\in V_k(F),\ F\subset\partial K,
\end{aligned}
\]
and with the moments \eqref{SigmaKDoFs2} set to zero.
Face moments are prescribed separately on each element side, using its
outward normal. Testing with the polynomial strain and projected jump gives
\[
\sum_K(\varepsilon_K^w(\boldsymbol v),\boldsymbol\sigma)_K
=\|\boldsymbol v\|_{1,h}^2.
\]
On the rescaled element, unisolvence, mesh regularity, and the uniform
basis bound in Assumption~\ref{meshassumption} give
\[
\begin{aligned}
\|\boldsymbol\sigma\|_K
&\lesssim\|\varepsilon(\boldsymbol\Pi_K^o\boldsymbol v^{\operatorname{div}})\|_K
+h_K^{-1/2}\|\boldsymbol\Pi_{\partial K}\mathcal J_{\partial K}(\boldsymbol v)\|_{\partial K},\\
\|\boldsymbol\sigma\|_{0,h}&\lesssim\|\boldsymbol v\|_{1,h}.
\end{aligned}
\]
Substitute this tensor into the supremum in \eqref{eq:discinfsup}.
The result follows for $\boldsymbol v\ne0$; the zero case is immediate.
\end{proof}}



\begin{lemma}\label{lem:coercive-strain}
For every $\boldsymbol v_h\in\boldsymbol V_h$,
\[
\|\boldsymbol v_h\|_{1,h}\lesssim\|\varepsilon_h^w(\boldsymbol v_h)\|.
\]
\end{lemma}
\begin{proof}
Apply Lemma~\ref{lem:discinf-sup} and the Cauchy--Schwarz inequality.
\end{proof}

\section{Numerical scheme}\label{sec:scheme}
Let
\[
S_h:=\{q_h\in L_0^2(\Omega):q_h|_K\in\mathbb P_{k-1}(K)
\text{ for all }K\in\mathcal K_h\}.
\]
The discrete formulation of \eqref{eq:stokes} seeks
$\boldsymbol u_h=\{\boldsymbol u_h^{\operatorname{div}},\boldsymbol u_h^t\}
\in\boldsymbol V_h$ and $p_h\in S_h$ such that
\begin{equation}\label{eq:scheme}
\left\{\begin{aligned}
a_h(\boldsymbol u_h,\boldsymbol v_h)+b(\boldsymbol v_h,p_h)
&=(\boldsymbol f,\boldsymbol\Pi_h^o\boldsymbol v_h^{\operatorname{div}})
&&\forall\boldsymbol v_h\in\boldsymbol V_h,\\
b(\boldsymbol u_h,q_h)&=0&&\forall q_h\in S_h,
\end{aligned}\right.
\end{equation}
where
\[
a_h(\boldsymbol u_h,\boldsymbol v_h):=
(\varepsilon_h^w(\boldsymbol u_h),\varepsilon_h^w(\boldsymbol v_h)),
\qquad
b(\boldsymbol v_h,p_h):=-(\operatorname{div}\boldsymbol v_h^{\operatorname{div}},p_h).
\]

\begin{lemma}
For every
\(
\boldsymbol{v}_h=\{\boldsymbol{v}_h^{\operatorname{div}},\boldsymbol{v}_h^t\}\in \boldsymbol{V}_h,
\)
there holds
\[
\|\varepsilon_h^w(\boldsymbol{v}_h)\|^2
\lesssim
\|\boldsymbol{v}_h\|_{1,h}^2.
\]
\end{lemma}

\begin{proof}
Set $\boldsymbol\sigma=\varepsilon_K^w(\boldsymbol v_h)$ in
\eqref{equ:weakstrain1}. The Cauchy--Schwarz and polynomial trace inequalities give
\[
\begin{aligned}
\|\varepsilon_K^w(\boldsymbol v_h)\|_K^2
\lesssim\Bigl(&\|\varepsilon(\boldsymbol\Pi_K^o\boldsymbol v_h^{\operatorname{div}})\|_K\\
&+h_K^{-1/2}\|\boldsymbol\Pi_{\partial K}\mathcal J_{\partial K}(\boldsymbol v_h)\|_{\partial K}\Bigr)
\|\varepsilon_K^w(\boldsymbol v_h)\|_K.
\end{aligned}
\]
Divide by the strain norm when it is nonzero, square, and sum over $K$.
\end{proof}


\begin{lemma}\label{lem:10}
For $\boldsymbol u_h,\boldsymbol v_h\in\boldsymbol V_h$,
\begin{align*}
|a_h(\boldsymbol{u}_h, \boldsymbol{v}_h)| &\lesssim \|\boldsymbol{u}_h\|_{1,h} \|\boldsymbol{v}_h\|_{1,h}, \\
a_h(\boldsymbol{u}_h, \boldsymbol{u}_h) &\gtrsim  \|\boldsymbol{u}_h\|_{1,h}^2.
\end{align*}
\end{lemma}




\begin{lemma}\label{lem:11}
For $\boldsymbol v_h\in\boldsymbol V_h$ and $q_h\in S_h$,
\[
|b(\boldsymbol v_h,q_h)|\lesssim\|\boldsymbol v_h\|_{1,h}\|q_h\|.
\]
\end{lemma}
\AI{\begin{proof}
Write $\boldsymbol w=\boldsymbol\Pi_K^o\boldsymbol v_h^{\operatorname{div}}$.
For $q\in\mathbb P_{k-1}(K)$, the volume projection and integration by parts give
\[
(\operatorname{div}\boldsymbol v_h^{\operatorname{div}},q)_K
=(\operatorname{div}\boldsymbol w,q)_K
+\langle\mathcal J_{\partial K}^{\boldsymbol n}(\boldsymbol v_h)\cdot\boldsymbol n,q\rangle_{\partial K}.
\]
Use $q=\operatorname{div}\boldsymbol v_h^{\operatorname{div}}$, the polynomial
trace inequality, and $|\operatorname{div}\boldsymbol w|\leq
\sqrt d\,|\varepsilon(\boldsymbol w)|$. This gives
$\|\operatorname{div}\boldsymbol v_h^{\operatorname{div}}\|_K
\lesssim\|\boldsymbol v_h\|_{1,h,K}$.
Sum over $K$ and apply the Cauchy--Schwarz inequality.
\end{proof}}

\begin{lemma}\label{lem:12}
For every $p_h\in S_h$, there is a $\boldsymbol v_h\in\boldsymbol V_h$ with
\[
\operatorname{div}\boldsymbol v_h^{\operatorname{div}}=p_h,
\qquad\|\boldsymbol v_h\|_{1,h}\lesssim\|p_h\|.
\]
\end{lemma}

\subsection{Interpolation}
For $\boldsymbol v\in\boldsymbol H^1(K)$, define
$\boldsymbol I_K^{\operatorname{div}}\boldsymbol v\in\boldsymbol V_K$
by matching its Type I--III DoFs with those of $\boldsymbol v$.
\AI{Integration by parts with $q\in\mathbb P_{k-1}(K)$ gives the commuting identity
\begin{equation}\label{eq:commuting}
\operatorname{div}\boldsymbol I_K^{\operatorname{div}}\boldsymbol v
=Q_h(\operatorname{div}\boldsymbol v)|_K,
\end{equation}
where $Q_h$ is the elementwise $L^2$ projection onto $\mathbb P_{k-1}$;
for mean-zero functions its range is $S_h$.
The interpolant reproduces $\mathbb P_k(K;\mathbb R^d)$.
Scaled stability of its DoFs gives
\begin{equation}\label{eq:Istability}
\|\boldsymbol I_K^{\operatorname{div}}\boldsymbol v\|_K
\lesssim\|\boldsymbol v\|_K+h_K|\boldsymbol v|_{1,K}.
\end{equation}
Indeed, the scaled face moments are bounded by
$h_K^{1/2}\|\boldsymbol v\cdot\boldsymbol n\|_{\partial K}$,
and the normalized volume moments by $\|\boldsymbol v\|_K$.
The trace inequality and uniform norm equivalence for the unisolvent DoFs
prove \eqref{eq:Istability}. These arguments are the usual face-VEM
interpolation estimates; see \cite{BBMR2016,BMM2022} for the underlying
spaces and related stability results. Subtracting a degree-$k$ polynomial
approximant in \eqref{eq:Istability} yields
\begin{equation}\label{eq:Iapprox}
\|\boldsymbol v-\boldsymbol I_K^{\operatorname{div}}\boldsymbol v\|_K
\lesssim h_K^{k+1}|\boldsymbol v|_{k+1,K}.
\end{equation}}

\AI{\begin{proof}[Proof of Lemma~\ref{lem:12}]
The continuous divergence right inverse \cite{BBF2013} provides
$\boldsymbol v\in\boldsymbol H_0^1(\Omega)$ with
$\operatorname{div}\boldsymbol v=p_h$ and
$|\boldsymbol v|_1\lesssim\|p_h\|$.
Set $\boldsymbol v_h^{\operatorname{div}}|_K=
\boldsymbol I_K^{\operatorname{div}}\boldsymbol v$ and
$\boldsymbol v_h^t|_F=\boldsymbol\Pi_{\partial K}^t\boldsymbol v$.
The traces are single-valued and zero on the boundary, and
\eqref{eq:commuting} gives the required divergence.
Let $\boldsymbol w=\boldsymbol\Pi_K^o\boldsymbol I_K^{\operatorname{div}}\boldsymbol v$.
Subtracting the mean of $\boldsymbol v$ in \eqref{eq:Istability}, and using
polynomial inverse estimates, gives
\[
\|\nabla\boldsymbol w\|_K\lesssim|\boldsymbol v|_{1,K},
\qquad
\|\boldsymbol v-\boldsymbol w\|_K\lesssim h_K|\boldsymbol v|_{1,K}.
\]
The normal interpolant trace is the face $L^2$ projection of
$\boldsymbol v\cdot\boldsymbol n$. Thus both parts of the projected boundary
jump are bounded by $\|\boldsymbol v-\boldsymbol w\|_{\partial K}$.
The trace inequality now gives
$\|\boldsymbol v_h\|_{1,h,K}\lesssim|\boldsymbol v|_{1,K}$.
Sum over $K$ to conclude.
\end{proof}
Since $b$ has a minus sign, use the negative of this lifting to obtain
\begin{equation}\label{eq:pressure-infsup}
\|q_h\|\lesssim\sup_{0\ne\boldsymbol v_h\in\boldsymbol V_h}
\frac{|b(\boldsymbol v_h,q_h)|}{\|\boldsymbol v_h\|_{1,h}}
\qquad(q_h\in S_h).
\end{equation}
Together with Lemmas~\ref{lem:10} and \ref{lem:11}, this proves that
\eqref{eq:scheme} has a unique solution \cite{BBF2013}.
Also, $\operatorname{div}\boldsymbol u_h^{\operatorname{div}}\in S_h$;
testing the second equation with this divergence shows that it vanishes
on every element.}

\subsection{Error equations}
Assume $(\boldsymbol u,p)\in
(\boldsymbol H^{k+1}(\Omega)\cap\boldsymbol H_0^1(\Omega))\times
(H^k(\Omega)\cap L_0^2(\Omega))$ solves \eqref{eq:stokes}.
Set
\[
\boldsymbol u_I^{\operatorname{div}}|_K=\boldsymbol I_K^{\operatorname{div}}\boldsymbol u,
\qquad \boldsymbol u_I^t|_F=\boldsymbol\Pi_{\partial K}^t\boldsymbol u,
\qquad p_I=Q_hp.
\]
\AI{For later use, \eqref{eq:Iapprox}, polynomial approximation, and inverse
estimates imply
\begin{equation}\label{eq:projected-approx}
\begin{aligned}
&\|\boldsymbol u-\boldsymbol\Pi_K^o\boldsymbol I_K^{\operatorname{div}}\boldsymbol u\|_K
+h_K|\boldsymbol u-\boldsymbol\Pi_K^o\boldsymbol I_K^{\operatorname{div}}\boldsymbol u|_{1,K}\\
&\quad+h_K^{1/2}\|\boldsymbol u-\boldsymbol\Pi_K^o\boldsymbol I_K^{\operatorname{div}}\boldsymbol u\|_{\partial K}
\lesssim h_K^{k+1}|\boldsymbol u|_{k+1,K}.
\end{aligned}
\end{equation}
To see this, subtract a degree-$k$ polynomial approximant and bound the
remaining polynomial by its $L^2$ norm using inverse and trace estimates.
The same argument applied to the strain gives
\begin{equation}\label{eq:strain-approx}
\begin{aligned}
&\|\varepsilon(\boldsymbol u)-\varepsilon(\boldsymbol\Pi_K^o
\boldsymbol I_K^{\operatorname{div}}\boldsymbol u)\|_K\\
&\quad+h_K^{1/2}\|\varepsilon(\boldsymbol u)-\varepsilon(\boldsymbol\Pi_K^o
\boldsymbol I_K^{\operatorname{div}}\boldsymbol u)\|_{\partial K}
\lesssim h_K^k|\boldsymbol u|_{k+1,K}.
\end{aligned}
\end{equation}
Since the normal interpolant trace and the tangential interpolant are face
projections, \eqref{eq:projected-approx} also gives
\begin{equation}\label{eq:interpolant-jump}
h_K^{-1/2}\|\boldsymbol\Pi_{\partial K}\mathcal J_{\partial K}(\boldsymbol u_I)\|_{\partial K}
\lesssim h_K^k|\boldsymbol u|_{k+1,K}.
\end{equation}}

\AI{Set $e_h=\boldsymbol u_h-\boldsymbol u_I$ and $\xi_h=p_h-p_I$.
The error equation is
\begin{equation}\label{eq:erroreq}
\begin{aligned}
&a_h(e_h,\boldsymbol v_h)+b(\boldsymbol v_h,\xi_h)
=\mathcal R(\boldsymbol v_h):=\sum_{K\in\mathcal K_h}\Bigl[\\
&\quad(\varepsilon(\boldsymbol u)-\varepsilon(\boldsymbol\Pi_K^o\boldsymbol u_I^{\operatorname{div}}),
 \varepsilon(\boldsymbol\Pi_K^o\boldsymbol v_h^{\operatorname{div}}))_K\\
&\quad+\langle\mathcal J_{\partial K}(\boldsymbol v_h),
 [\varepsilon(\boldsymbol u)-\varepsilon(\boldsymbol\Pi_K^o\boldsymbol u_I^{\operatorname{div}})]\boldsymbol n\rangle_{\partial K}\\
&\quad+\langle Q_hp-p,
 (\boldsymbol v_h^{\operatorname{div}}-\boldsymbol\Pi_K^o\boldsymbol v_h^{\operatorname{div}})\cdot\boldsymbol n\rangle_{\partial K}\\
&\quad-\langle\boldsymbol\Pi_{\partial K}\mathcal J_{\partial K}(\boldsymbol u_I),
 \varepsilon_K^w(\boldsymbol v_h)\boldsymbol n\rangle_{\partial K}\Bigr]
\qquad\forall\boldsymbol v_h\in\boldsymbol V_h.
\end{aligned}
\end{equation}
To derive it, test \eqref{eq:stokes} with
$\boldsymbol\Pi_K^o\boldsymbol v_h^{\operatorname{div}}$, integrate by parts,
and sum over $K$. The exact traction paired with
$(\boldsymbol v_h^{\operatorname{div}}\cdot\boldsymbol n)\boldsymbol n+\boldsymbol v_h^t$
cancels on interior faces and vanishes on the boundary.
For the pressure terms, the volume projection preserves moments against
$\nabla(Q_hp)$, and
\[
(\operatorname{div}\boldsymbol\Pi_K^o\boldsymbol v_h^{\operatorname{div}},p-Q_hp)_K=0.
\]
For the strain terms, use \eqref{equ:weakstrain1} twice and recall that
$\varepsilon(\boldsymbol\Pi_K^o\boldsymbol u_I^{\operatorname{div}})
\in\mathbb P_{k-1}(K;\mathbb S)\subset\Sigma_h^{-1}(K)$.
Its normal trace belongs to $V_k(F)$, so it pairs equally with the projected
and full jumps.
By \eqref{eq:commuting}, $\operatorname{div}\boldsymbol u_I^{\operatorname{div}}=0$.
Hence
\begin{equation}\label{eq:divergenceerror}
b(e_h,q_h)=0\qquad\forall q_h\in S_h.
\end{equation}}

\AI{We estimate the four terms in \eqref{eq:erroreq} in order.
The volume strain term is bounded by \eqref{eq:strain-approx} and the norm
of $\boldsymbol v_h$. The boundary strain term uses the same approximation
estimate and the full-jump bound \eqref{eq:fulljump}.
For the pressure term, the local projection estimate is
\[
h_K^{1/2}\|p-Q_hp\|_{\partial K}\lesssim h_K^k|p|_{k,K};
\]
its other factor is the normal part of the jump, which the discrete norm
controls directly. Finally, the polynomial trace inequality and the weak-strain
upper bound give
\[
\begin{aligned}
&|\langle\boldsymbol\Pi_{\partial K}\mathcal J_{\partial K}(\boldsymbol u_I),
\varepsilon_K^w(\boldsymbol v_h)\boldsymbol n\rangle_{\partial K}|\\
&\qquad\lesssim h_K^{-1/2}
\|\boldsymbol\Pi_{\partial K}\mathcal J_{\partial K}(\boldsymbol u_I)\|_{\partial K}
\|\boldsymbol v_h\|_{1,h,K}.
\end{aligned}
\]
Use \eqref{eq:interpolant-jump}, then sum the estimates over $K$.
The Cauchy--Schwarz inequality yields
\begin{equation}\label{eq:residual-bound}
|\mathcal R(\boldsymbol v_h)|\lesssim
h^k(\|\boldsymbol u\|_{k+1}+\|p\|_k)\|\boldsymbol v_h\|_{1,h}.
\end{equation}}

\subsection{Error estimates}
Taking $\boldsymbol v_h=e_h$ in \eqref{eq:erroreq} and
$q_h=\xi_h$ in \eqref{eq:divergenceerror}, coercivity gives
\[
\|e_h\|_{1,h}^2\lesssim a_h(e_h,e_h)=\mathcal R(e_h).
\]
Therefore, \eqref{eq:residual-bound} implies
\[
\|e_h\|_{1,h}\lesssim h^k(\|\boldsymbol u\|_{k+1}+\|p\|_k).
\]
\AI{For the pressure, \eqref{eq:pressure-infsup} and \eqref{eq:erroreq} give
\[
\begin{aligned}
\|\xi_h\|&\lesssim\sup_{0\ne\boldsymbol v_h\in\boldsymbol V_h}
\frac{|\mathcal R(\boldsymbol v_h)-a_h(e_h,\boldsymbol v_h)|}{\|\boldsymbol v_h\|_{1,h}}\\
&\lesssim\|e_h\|_{1,h}+h^k(\|\boldsymbol u\|_{k+1}+\|p\|_k).
\end{aligned}
\]
Thus $\|\xi_h\|$ has the same bound.}

\AI{In the velocity error norm, identify the smooth velocity with the pair
$\{\boldsymbol u,\boldsymbol\Pi_{\partial K}^t\boldsymbol u\}$, and use the same
norm formula with the $L^2$ volume projector.
Polynomial approximation and \eqref{eq:Iapprox} give
$\|\boldsymbol u-\boldsymbol u_I\|_{1,h}\lesssim h^k\|\boldsymbol u\|_{k+1}$.
The triangle inequality and $\|p-Q_hp\|\lesssim h^k\|p\|_k$ then yield
\begin{equation}\label{eq:final-error}
\|\boldsymbol u-\boldsymbol u_h\|_{1,h}+\|p-p_h\|
\lesssim h^k(\|\boldsymbol u\|_{k+1}+\|p\|_k).
\end{equation}}

% !TEX root = SEVEM_Stokes.tex
\section{Numerical results}\label{sec:numerics}

\AI{For $k=1$, the velocity uses $\boldsymbol V_h$ and the pressure is
piecewise constant on the primal mesh $\mathcal K_h$.
The coefficient of the symmetric strain in \eqref{eq:stokes} is one,
which corresponds to $2\nu=1$ in the usual viscosity convention.
The weak strain \eqref{equ:weakstrain1} has the equivalent form
\[
(\varepsilon_K^w(\boldsymbol v),\boldsymbol\sigma)_K
=-(\boldsymbol v^{\operatorname{div}},\operatorname{div}_{w,K}\boldsymbol\sigma)_K
+\langle\boldsymbol v^{\operatorname{div}}\cdot\boldsymbol n,
 (\boldsymbol\sigma\boldsymbol n)\cdot\boldsymbol n\rangle_{\partial K}
+\langle\boldsymbol v^t,\Pi_F(\boldsymbol\sigma\boldsymbol n)\rangle_{\partial K},
\]
where $\boldsymbol\sigma\in\Sigma_h^{-1}(K)$ with $k=1$, and
$\operatorname{div}_{w,K}\boldsymbol\sigma\in\mathbb P_1(K;\mathbb R^d)$ satisfies
\[
(\operatorname{div}_{w,K}\boldsymbol\sigma,\boldsymbol q)_K
=-(\boldsymbol\sigma,\nabla\boldsymbol q)_K
+\langle\boldsymbol\sigma\boldsymbol n,\boldsymbol q\rangle_{\partial K}
\qquad\forall\boldsymbol q\in\mathbb P_1(K;\mathbb R^d).
\]
This definition includes the internal stress jumps through integration by
parts on the cones. Its equivalence follows from the volume projection
and the symmetry of $\boldsymbol\sigma$.}

The reported \texttt{sincosdata} test is posed on the unit square.
Tables~\ref{tab:poly}--\ref{tab:rect} give the supplied results.
\AI{The energy error is $\|\boldsymbol u_I-\boldsymbol u_h\|_{A_h}
:=a_h(\boldsymbol u_I-\boldsymbol u_h,\boldsymbol u_I-\boldsymbol u_h)^{1/2}$.
The energy and total pressure errors approach first-order convergence on
these meshes. Faster convergence in some other columns is an empirical
observation; it is not proved by \eqref{eq:final-error}.}

\AI{\begin{remark}[Numerical data to confirm]
The supplied files do not specify the exact velocity, pressure, force,
boundary data, viscosity used in the computations, or the definitions of
$\|\cdot\|_{0,h}$ and the reported maximum norms.
The implemented strain space must also be checked against
\eqref{eq:strainspace}; the original numerical description used a different
index. The table values are retained pending these checks.
The rectangular meshes do not satisfy Assumption~\ref{meshassumption}:
their face-tangential tensors span only diagonal symmetric matrices.
Table~\ref{tab:rect} therefore reports a test outside the proved mesh class.
The three mesh figure files were not supplied.
\end{remark}}

% ===== Table 1 =====
\begin{table}[htbp]
  \centering \footnotesize
  \caption{Results on polygonal meshes obtained from a triangle-mesh dual.}\label{tab:poly}
  \setlength{\tabcolsep}{3pt}\resizebox{\textwidth}{!}{\begin{tabular}{cccccccc}
    \toprule
    DoFs & $h$ & $\|\boldsymbol u_I-\boldsymbol u_h\|_{A_h}$ & $\|\boldsymbol u_I-\boldsymbol u_h\|_{0,h}$ & $\|\boldsymbol u_I-\boldsymbol u_h\|_{\infty}$ & $\|p_I-p_h\|$ & $\|p-p_h\|$ & $\|p_I-p_h\|_{\text{max}}$ \\
    \midrule
    314   & 1.43e-01 & 8.767e-01 & 2.588e-02 & 7.996e-02 & 2.071e-01 & 2.454e-01 & 5.002e-01 \\
    978   & 7.78e-02 & 4.012e-01 & 9.135e-03 & 2.645e-02 & 6.612e-02 & 1.007e-01 & 1.870e-01 \\
    3362  & 4.12e-02 & 1.884e-01 & 2.712e-03 & 7.111e-03 & 1.936e-02 & 4.463e-02 & 7.481e-02 \\
    12354 & 2.13e-02 & 9.119e-02 & 7.279e-04 & 1.817e-03 & 5.869e-03 & 2.145e-02 & 3.068e-02 \\
    47234 & 1.08e-02 & 4.490e-02 & 1.875e-04 & 4.573e-04 & 1.874e-03 & 1.061e-02 & 1.335e-02 \\
    \bottomrule
  \end{tabular}}
\end{table}

% ===== Table 2 =====
\begin{table}[htbp]
  \centering\footnotesize
  \caption{Results on triangular meshes.}\label{tab:tri}
  \setlength{\tabcolsep}{3pt}\resizebox{\textwidth}{!}{\begin{tabular}{cccccccc}
    \toprule
    DoFs & $h$ & $\|\boldsymbol u_I-\boldsymbol u_h\|_{A_h}$ & $\|\boldsymbol u_I-\boldsymbol u_h\|_{0,h}$ & $\|\boldsymbol u_I-\boldsymbol u_h\|_{\infty}$ & $\|p_I-p_h\|$ & $\|p-p_h\|$ & $\|p_I-p_h\|_{\text{max}}$ \\
    \midrule
    232   & 2.50e-01 & 4.289e-01 & 1.463e-02 & 3.022e-02 & 5.357e-01 & 5.509e-01 & 8.008e-01 \\
    880   & 1.25e-01 & 1.949e-01 & 2.959e-03 & 6.035e-03 & 2.867e-01 & 2.940e-01 & 4.262e-01 \\
    3424  & 6.25e-02 & 9.410e-02 & 6.990e-04 & 1.479e-03 & 1.450e-01 & 1.486e-01 & 2.115e-01 \\
    13504 & 3.12e-02 & 4.659e-02 & 1.727e-04 & 4.138e-04 & 7.265e-02 & 7.446e-02 & 1.045e-01 \\
    53632 & 1.56e-02 & 2.323e-02 & 4.308e-05 & 1.098e-04 & 3.634e-02 & 3.725e-02 & 5.187e-02 \\
    \bottomrule
  \end{tabular}}
\end{table}

% ===== Table 3 =====
\begin{table}[htbp]
  \centering\footnotesize
  \caption{Results on rectangular meshes, outside Assumption~\ref{meshassumption}.}\label{tab:rect}
  \setlength{\tabcolsep}{3pt}\resizebox{\textwidth}{!}{\begin{tabular}{cccccccc}
    \toprule
    DoFs & $h$ & $\|\boldsymbol u_I-\boldsymbol u_h\|_{A_h}$ & $\|\boldsymbol u_I-\boldsymbol u_h\|_{0,h}$ & $\|\boldsymbol u_I-\boldsymbol u_h\|_{\infty}$ & $\|p_I-p_h\|$ & $\|p-p_h\|$ & $\|p_I-p_h\|_{\text{max}}$ \\
    \midrule
    152    & 2.50e-01 & 6.193e-01 & 4.377e-02 & 7.423e-02 & 1.557e-02 & 1.574e-01 & 2.227e-02 \\
    560    & 1.25e-01 & 2.438e-01 & 1.277e-02 & 1.582e-02 & 9.016e-03 & 8.021e-02 & 1.778e-02 \\
    2144   & 6.25e-02 & 1.083e-01 & 3.348e-03 & 4.464e-03 & 2.688e-03 & 4.011e-02 & 6.658e-03 \\
    8384   & 3.12e-02 & 5.211e-02 & 8.477e-04 & 1.148e-03 & 6.957e-04 & 2.004e-02 & 1.943e-03 \\
    33152  & 1.56e-02 & 2.579e-02 & 2.126e-04 & 2.889e-04 & 1.752e-04 & 1.002e-02 & 5.205e-04 \\
    \bottomrule
  \end{tabular}}
\end{table}

% The figure is included only when all three supplied assets are available.
\IfFileExists{fig/poly.pdf}{%
  \IfFileExists{fig/tri.pdf}{%
    \IfFileExists{fig/rect.pdf}{%
      \begin{figure}[htbp]
      \centering
      \includegraphics[width=.31\textwidth]{fig/poly.pdf}\hfill
      \includegraphics[width=.31\textwidth]{fig/tri.pdf}\hfill
      \includegraphics[width=.31\textwidth]{fig/rect.pdf}
      \caption{Polygonal, triangular, and rectangular test meshes.}\label{fig:meshes}
      \end{figure}
    }{}%
  }{}%
}{}%



\clearpage
\AI{\begin{thebibliography}{8}
\bibitem{BBCMMR2013}
L. Beir\~ao da Veiga, F. Brezzi, A. Cangiani, G. Manzini, L. D. Marini, and A. Russo,
\emph{Basic principles of virtual element methods},
Math. Models Methods Appl. Sci., 23 (2013), pp. 199--214,
\url{https://doi.org/10.1142/S0218202512500492}.

\bibitem{BBMR2016}
L. Beir\~ao da Veiga, F. Brezzi, L. D. Marini, and A. Russo,
\emph{$H(\mathrm{div})$ and $H(\mathrm{curl})$-conforming virtual element methods},
Numer. Math., 133 (2016), pp. 303--332,
\url{https://doi.org/10.1007/s00211-015-0746-1}.

\bibitem{BMM2022}
L. Beir\~ao da Veiga, L. Mascotto, and J. Meng,
\emph{Interpolation and stability estimates for edge and face virtual elements of general order},
Math. Models Methods Appl. Sci., 32 (2022), pp. 1589--1631,
\url{https://doi.org/10.1142/S0218202522500373}.

\bibitem{BBF2013}
D. Boffi, F. Brezzi, and M. Fortin,
\emph{Mixed Finite Element Methods and Applications},
Springer Series in Computational Mathematics, vol. 44, Springer, Heidelberg, 2013,
\url{https://doi.org/10.1007/978-3-642-36519-5}.

\bibitem{Brenner2004}
S. C. Brenner,
\emph{Korn's inequalities for piecewise $H^1$ vector fields},
Math. Comp., 73 (2004), pp. 1067--1087,
\url{https://doi.org/10.1090/S0025-5718-03-01579-5}.

\bibitem{ChenWang2019}
L. Chen and F. Wang,
\emph{A divergence free weak virtual element method for the Stokes problem on polytopal meshes},
J. Sci. Comput., 78 (2019), pp. 864--886,
\url{https://doi.org/10.1007/s10915-018-0796-5}.

\bibitem{ChenHuangWei2024}
C. Chen, X. Huang, and H. Wei,
\emph{Virtual element methods without extrinsic stabilization},
SIAM J. Numer. Anal., 62 (2024), pp. 567--591,
\url{https://doi.org/10.1137/22M1504196}.

\bibitem{MuYeZhang2021}
L. Mu, X. Ye, and S. Zhang,
\emph{A stabilizer-free, pressure-robust, and superconvergence weak Galerkin
finite element method for the Stokes equations on polytopal mesh},
SIAM J. Sci. Comput., 43 (2021), pp. A2614--A2637,
\url{https://doi.org/10.1137/20M1380405}.
\end{thebibliography}}

\end{document} 
